本书旨在向读者介绍偏微分方程的数学理论. 但是, 如果暗示偏微分方程存在一种与常微分方程理论或单复变函数理论意义相同的"理论", 就会产生误导. 偏微分方程比上述两个学科广泛得多 (它将二者都包含为特殊情形), 而且发展得还不够完善. 不过, 尽管一个随意的观察者可能认为, 这一学科不过是用于处理不同类型问题的一袋互不相关的技巧, 实际上仍有某些主题贯穿始终.

为了说明这些主题, 我们采用两种途径. 第一种是提出一组在许多偏微分方程问题中都会出现的问题 (存在性、多解性等). 为了举例说明处理这些问题的不同方法, 我们考察常微分方程理论、高等微积分和复变函数理论中的一些结果 (假定读者对这些内容已有一定了解). 第二种途径是以非常初等的方式考察三个偏微分方程 (Laplace 方程、热方程和波动方程); 同样, 对大多数读者而言, 这可能只是复习. 我们将看到, 即使最初等的方法, 也预示着本书后面各章中的更深刻结果.
